% !TeX root = ../../tesis.tex

% =============================================================================
\subsubsection{Capas del sistema de información geográfica}
\label{subsubsec:tg-capas}
% =============================================================================

El módulo de datos mantiene diez capas \termino{GeoPackage} en
\texttt{data/processed/}, todas proyectadas en \textsc{epsg}:32614 (UTM
zona~14N), referencia cartográfica adecuada para la extensión de la \zmvm{}. Las
capas se organizan en tres categorías temáticas, ilustradas en la Figura~
\ref{fig:tg-capas}: (i)~\textbf{contexto territorial}, con los límites
político-administrativos de la \cdmx{}, del Estado de México, el polígono de
estudio y las variables de medio físico natural; (ii)~
\textbf{red e infraestructura}, con las geometrías de los sistemas de
transporte, la infraestructura vial y el trazado del Periférico; y (iii)~
\textbf{indicadores del modelo VFT}, con las capas derivadas del cálculo de
cobertura espacial, fuerza capilar y factor de desviación descritos en el
Capítulo~ \ref{ch:resultados-vft}.

\begin{figure}[htbp]
\centering
% Diagrama: Capas GeoPackage de Transport-gis organizadas por categoría
% Usado en: 4-Capas-Codigo/4-3-Visualizacion.tex  (fig:tg-capas)
% Fix 2026-08-28: node distance aumentado de 0.12cm a 0.28cm para dar respiración al texto
\begin{tikzpicture}[
  hdr1/.style = {rectangle, rounded corners=3pt,
                 draw=gray!60, fill=gray!55,
                 minimum width=4.2cm, minimum height=0.80cm,
                 align=center, font=\small\bfseries\color{white}},
  lyr1/.style = {rectangle, rounded corners=2pt,
                 draw=gray!45, fill=gray!8,
                 minimum width=4.2cm, minimum height=0.72cm,
                 align=center, font=\footnotesize},
  hdr2/.style = {rectangle, rounded corners=3pt,
                 draw=unamAzul, fill=unamAzul!70,
                 minimum width=4.2cm, minimum height=0.80cm,
                 align=center, font=\small\bfseries\color{white}},
  lyr2/.style = {rectangle, rounded corners=2pt,
                 draw=unamAzul!50, fill=unamAzul!8,
                 minimum width=4.2cm, minimum height=0.72cm,
                 align=center, font=\footnotesize},
  hdr3/.style = {rectangle, rounded corners=3pt,
                 draw=unamOro!80, fill=unamOro!75,
                 minimum width=4.2cm, minimum height=0.80cm,
                 align=center, font=\small\bfseries\color{white}},
  lyr3/.style = {rectangle, rounded corners=2pt,
                 draw=unamOro!60, fill=unamOro!10,
                 minimum width=4.2cm, minimum height=0.72cm,
                 align=center, font=\footnotesize},
  node distance = 0.28cm
]
%% — Columna 1: Contexto territorial —
\node[hdr1] (h1) at (0,0) {Contexto territorial};
\node[lyr1, below=of h1] (c1a) {\texttt{LimitesPoliticos.gpkg}};
\node[lyr1, below=of c1a] (c1b) {\texttt{PoligonosEdoMex.gpkg}};
\node[lyr1, below=of c1b] (c1c) {\texttt{PoligonoEstudio1.gpkg}};
\node[lyr1, below=of c1c] (c1d) {\texttt{MedioFisicoNatural.gpkg}};

%% — Columna 2: Red e infraestructura —
\node[hdr2] (h2) at (4.8,0) {Red e infraestructura};
\node[lyr2, below=of h2] (c2a) {\texttt{Transporte.gpkg}};
\node[lyr2, below=of c2a] (c2b) {\texttt{InfraEstructura.gpkg}};
\node[lyr2, below=of c2b] (c2c) {\texttt{Vialidades.gpkg}};

%% — Columna 3: Indicadores del modelo —
\node[hdr3] (h3) at (9.6,0) {Indicadores del modelo};
\node[lyr3, below=of h3] (c3a) {\texttt{Cobertura800m.gpkg}};
\node[lyr3, below=of c3a] (c3b) {\texttt{FuerzaCapilar.gpkg}};
\node[lyr3, below=of c3b] (c3c) {\texttt{FactorDesviacion.gpkg}};
\end{tikzpicture}

\caption[Organización de las capas GeoPackage de Transport-gis]%
{Las diez capas \termino{GeoPackage} del directorio \texttt{data/processed/},
organizadas por
categoría temática. Todas están proyectadas en \textsc{epsg}:32614 (UTM
zona~14N) y versionadas con Git~\textsc{lfs}.}
\label{fig:tg-capas}
\end{figure}
